  \mysubsec{Exercício 1}
\bigline
\textbf{Seja $X$ um conjunto não-enumerável. Considere a seguinte coleção:
$$\mathcal F = 
\{A \subseteq X ; A \text{ ou } A^C \text{ é enumerável}\}.$$ 
  Mostre que $\mathcal F$ é uma $\sigma$-álgebra.}
\\

\begin{proof}[\textbf{Dem:}]
 Ora, visto que $\emptyset$ é enumerável, $\emptyset \in \mathcal F$. 
E, já que $\emptyset = X^C$, temos $X \in \mathcal F$. Mais ainda,
para qualquer $A \in \mathcal F$ temos que $A$ ou $A^C$ é enumerável. 
Em ambos os casos, podemos concluír que $A^C \in \mathcal F$. Desse modo, $\mathcal F$ 
satisfaz as duas primeiras condições de uma $\sigma$-álgebra de $X$.

Para verificar a terceira condição, seja 
$\{A_n\}_{n \in \mathbb N} \subset \mathcal F$. Note que 
podemos recair a dois casos, são eles:
\begin{enumerate}
  \item Existe $k \in \mathbb N$ tal que $A_k^C$ é enumerável;
  \item Para todo $n \in \mathbb N$ temos $A_n^C$ não-enumerável. 
\end{enumerate}

Assumindo (1), já que subconjuntos de conjuntos enumeráveis são 
enumeráveis, segue que $\bigcap A_n^C \subset A_k^C$ é enumerável, ou seja 
$\left(\bigcap A_n^C\right)^C = \bigcup A_n \in \mathcal F$. 

Caso contrário, assumindo (2), já que para todo $n \in \mathbb N$, 
temos $A_n^C$ não-enumerável e $A_n \in \mathcal F$, então todo $A_n$ é enumerável.
Desse modo, visto que a união enumerável de enumeráveis é enumerável, temos 
$\bigcup A_n$ enumerável, ou seja $\bigcup A_n \in \mathcal F$. Portanto 
$\mathcal F$ satisfaz as condições de uma $\sigma$-álgebra de $X$, 
como queríamos. 
\end{proof}
\bigline
\newpage
\mysubsec{Exercício 2}
\bigline
\textbf{Caracterize a $\sigma$-álgebra de Borel de $\mathbb R$.}
\noindent\textbf{} 
\begin{proof}[\textbf{Dem:}]
Considere os seguintes conjuntos:
\begin{align*}
  &E_1 = \{(a, b) ; a<b\}, &&E_2= \{[a, b]; a< b\},\\
  &E_3 = \{(a, b] ; a<b\}, &&E_4= \{[a, b); a< b\},\\
  &E_5 = \{(a, +\infty) ; a\in\mathbb R\}, &&E_6= \{(+\infty, a); a \in \mathbb R\}, \\
  &E_7 = \{[a, +\infty) ; a\in\mathbb R\}, &&E_8= \{(+\infty, a]; a \in \mathbb R\}.
\end{align*}
Vejamos que para cada $n \in \mathbb N_{\leq 8}$, $\mathcal B(\mathbb R, \tau)$
é gerada por $E_n$.

(i) ($\subseteq$) Note que $E_1 \subseteq \tau \subset \sigma(\tau)$, 
desse modo $\sigma(E_1) \subseteq \sigma(\tau)$.\\ 

(i)($\supseteq$) Por outro lado, visto que todo elemento de $\tau$ é um aberto
e, todo aberto pode ser escrito como a união enumerável de elementos de $E_1$, 
segue que $\tau \subseteq \sigma(E_1)$. Com isso, temos $\sigma(\tau) \subseteq \sigma(E_1)$.
Dito isso, segue que $E_1$ gera $\sigma(\tau)=\mathcal B (\mathbb R)$. 

(ii) ($\subseteq$) Ora, para cada $[a, b] \in E_2$, 
$$[a, b] = \bigcap\limits_{n \in \mathbb N} (a - 1/n, b + 1/n),$$
portanto $E_2 \subseteq \sigma(E_1)$, isto é, $\sigma(E_2) \subseteq \sigma(E_1)$. 

(ii) ($\supseteq$) Por outro lado, para cada $(a, b) \in E_1$,
$$(a, b) = \bigcup\limits_{n \ge 3} [a + (b-a)/n, b - (b-a)/n],$$
desse modo, $E_1 \subseteq \sigma(E_2)$, ou seja, $\sigma(E_1) \subseteq \sigma(E_2)$.\\
Sendo assim, $\sigma(E_1) = \sigma(E_2)$, isto é, $E_2$ gera $\mathcal B(\mathbb R)$.

(iii) e (iv) Já que cada elemento de $E_1$, $E_3$ e $E_4$ 
podem ser rescritos das seguintes formas: 
\begin{align*}
  &(a, b] = \bigcap\limits_{n \in \mathbb N} (a, b+1/n), 
  &&(a, b) = \bigcup\limits_{n \ge 2} (a, b-(b-a)/n],\\
  &[a, b) = \bigcap\limits_{n \in \mathbb N} (a-1/n, b), 
  &&(a, b) = \bigcup\limits_{n \ge 2} [a+(b-a)/n, b).
\end{align*}
Segue de argumentos similares a (ii) que $E_3$ e $E_4$ geram $\mathcal B(\mathbb R, \tau)$.

  (v), (vi), (vii) e (viii) Note que $\sigma(E_5) = \sigma(E_8)$, $\sigma(E_6) = \sigma(E_7)$ e 
  $\sigma(E_5) = \sigma(E_7)$ por argumentos semelhantes a (i). Daqui, visto que 
\begin{align*}
  &(a, +\infty) = \bigcup\limits_{n \in \mathbb N} (a, b+n), 
  &(a, b) = (a, +\infty) \cap \left(\bigcup\limits_{n \in \mathbb N} (b-1/n, +\infty) \right)^C
\end{align*}
  temos $\sigma(E_1) = \sigma(E_5)$. Com isso, para todo $n \in \mathbb N_{\le 8}$, $E_n$ gera 
  $\mathcal B(\mathbb R, \tau)$, como queríamos.
\end{proof}
\bigline

\newpage
\mysubsec{Exercício 3}
\bigline 
\textbf{Para $E \in \mathcal B(\mathbb R)$, considere $E_1 = E \cup \{+\infty\}$, 
$E_2 = E \cup \{-\infty\}$ e $E_3 = E \cup \{-\infty, +\infty\}$. Prove que 
$\{E, E_1, E_2, E_3; E\in \mathcal B(\mathbb R)\}$ 
é uma $\sigma$-álgebra de $\overline{\mathbb R}$}. 
\begin{proof}[\textbf{Dem:}] Defina $\mathcal C = \{E, E_1, E_2, E_3; E
\in \mathcal B(\mathbb R)\}$,
visto que $\emptyset\in \mathcal B(\mathbb R)$, segue que $\emptyset \in \mathcal C$.
De maneira análoga, temos $\overline{\mathbb R} \in \mathcal C$. Agora, seja $A \in C$, 
\begin{itemize}
  \item Para  $A \in \mathcal B(\mathbb R)$, 
    segue que $\mathbb R\setminus A \in \mathcal B(\mathbb R)$, sendo assim,
    $\left(\mathbb R \setminus A \right)\cup \{-\infty, +\infty\} = A^C \in \mathcal C$.

  \item Para $A = E \cup \{+ \infty\}$,
    note que $A^C = X \cup \{- \infty\}$ com 
    $X \in \mathcal B (\mathbb R)$, sendo assim, $A^C \in \mathcal C$. 
  \item Para $A = E \cup \{- \infty\}$,
    segue de maneira análoga ao 
    caso anterior que $A^C \in \mathcal C$. 
  \item Para $A = E \cup \{- \infty, + \infty\}$, temos $A^C \in \mathcal B(\mathbb R)$, 
    portanto, $A^C \in \mathcal C$
\end{itemize}
Para verificar que $\mathcal C$ obedesce a terceira propriedade de $\sigma$-álgebra,
note que todo $A \in \mathcal C$ 
temos $A = B \cup D$, onde $B \in \mathcal B (\mathbb R)$ e 
$D \in \mathbb P (\{-\infty, +\infty\}) $. Portanto, dada qualquer 
sequência $\{A_k\}_{k \in \mathbb N}
\subset \mathcal C$, temos $\bigcup A_k = \bigcup \left(B_k \cup D_k\right)$ para algum 
$\{B_k\} \subset \mathcal B (\mathbb R)$ e 
  $\{D_k\}\subset \mathbb P (\{-\infty, +\infty\})$. Daqui, já que 
\begin{align*}
  \bigcup A_k &= \bigcup \left(B_k \cup D_k\right),\\
  &= \left(\bigcup B_m\right) \cup \left(\bigcup D_n\right), 
\end{align*}
e, $\bigcup D_n \in \mathbb P (\{-\infty, +\infty\}) \subset \mathcal C$, assim como 
  $\bigcup B_m \in \mathcal B(\mathbb R) \subset \mathcal C$, então 
  $\bigcup A_k \in \mathcal C$.
\end{proof}
\bigline
\newpage
\mysubsec{Exercício 4}
\bigline 
\textbf{Sejam $f \in M(X, \mathcal F)$ e $n \in \mathbb N$. Considere o truncamento de $f$, 
$f_n: X \rightarrow \mathbb R$ definido por
\begin{align*}
  f_n(x) =\begin{cases}
    -n, &\text{se } f(x) < -n,\\
    f(x), &\text{se } |f(x)| \le n,\\
    n, &\text{se } f(x)>n.
  \end{cases}
\end{align*}
Mostre que $f_n$ é mensurável para todo $n \in \mathbb N$ e que $\lim f_n(x) = f(x)$, 
para todo $x \in X$.}
\begin{proof}[\textbf{Dem:}] Fixado $n \in \mathbb N$, seja $\alpha \in \mathbb R$. 

\begin{itemize}
  \item Para $\alpha < -n$, temos $\{x \in X; f_n(x) > \alpha\}
  = X \in \mathcal F$;
  \item   Para $-n \le \alpha < n$, $f_n = f$, sendo assim, $\{x \in X; f_n(x) > \alpha\} = 
    \{x \in X; f(x) > \alpha\}$. Visto que $f$ é mensurável, 
    $\{x \in X; f_n(x) > \alpha\} \in \mathcal F$;
  \item Por fim, para $\alpha \ge n$, $\{x \in X; f(x) > \alpha\} = \emptyset \in \mathcal F$.
\end{itemize}
Portanto, para todo $n \in \mathbb N$ temos $f_n$ mensurável. 

Vejamos que $\lim f_n(x) = f(x)$.
Dado $x \in X$, como $f(x) \in \mathbb R$, existe $M \in \mathbb N$ tal que $-M\le f(x)\le M$ 
(e.g o menor inteiro maior que o valor absoluto de $f(x)$).
Desse modo, para inteiros $n \ge M$, temos $f_n(x) = f(x)$, 
ou seja, para todo $\epsilon > 0$, $|f_n(x) - f(x)| = 0 < \epsilon$, 
isto é, $\lim f_n (x) = f(x)$, como queríamos.
\end{proof}
\bigline
\newpage 
\mysubsec{Exercício 5}
\bigline 
\textbf{Mostre que a soma, o produto, o valor absoluto, e limites de funções mensuráveis
a valores complexos são mensuráveis.}
\begin{proof}[\textbf{Dem:}] Por definição, uma função a valores complexos é mensurável se, 
  e somente se, a parte real e a parte imaginária da função são funções mensuráveis. 
  Desse modo, sejam $f, g: X \rightarrow \mathbb C$ funções mensuráveis, tais que 
  $f_1, f_2, g_1, g_2: X \rightarrow \mathbb R$ satisfazem $f(x) = f_1(x)+if_2(x)$
  e $g(x)=g_1(x)+ig_2(x)$ para todo $x \in X$.
  \vspace{5mm}

\noindent
(Soma:) Ora,
\begin{align*}
  (f+g)(x) &= f(x) + g(x) = [f_1(x)+if_2(x)] + [g_1(x)+ig_2(x)],\\
  &= [f_1(x) + g_1(x)] + i[f_2(x) + g_2(x)].
\end{align*}
Visto que $f_1(x)$ e $g_1(x)$ são funções mensuráveis, temos $(f_1 + g_1)(x)$ mensurável. Analogamente temos 
  $(f_2+g_2)(x)$ mensurável. Já que a parte real e a parte imaginária de $(f+g)(x)$ é mensurável, temos 
  $(f+g)(x)$ mensurável, como queríamos.
  \vspace{5mm}

\noindent
  (Produto:) Já que, 
\begin{align*}
  (f\cdot g)(x) &= f(x)\cdot g(x) = [f_1(x)+if_2(x)]\cdot[g_1(x) + ig_2(x)], \\
  &= [(f_1\cdot g_1)(x) - (f_2\cdot g_2)(x)] + i[(f_1\cdot g_2)(x) + (f_2 \cdot g_1)(x)],\\
\end{align*}
  e temos $f_n,g_n: X \rightarrow \mathbb R$ mensuráveis para $n=1, 2$, segue que $(f_i\cdot g_j)(x)$ é 
  mensurável para todo par $i, j \in \{1, 2\}$. Sendo assim, $(f_1\cdot g_1)(x) - (f_2\cdot g_2)(x)$ e 
  $(f_1\cdot g_2)(x) + (f_2 \cdot g_1)(x)$ são também mensuráveis, ou seja $(f\cdot g)(x)$ é mensurável. 
  
  \vspace{5mm}
  \noindent
  (Valor Absoluto:) A demonstração será divida nas seguintes partes:\\

  Para $\alpha \ge 0$, 
  \begin{align*}
    \{x \in X;|f(x)|> \alpha\} = \{ x \in X;(f_1(x))^2 + (f_2(x))^2 > \alpha^2\},
  \end{align*}
  Já que $f_1, f_2$ são funções mensuráveis, a soma dos quadrados é mensurável, i.e.
  \begin{align*}
   \{x \in X;|f(x)|> \alpha\} = \{ x \in X;(f_1(x))^2 + (f_2(x))^2 > \alpha^2\}\in \mathcal F.
  \end{align*}

  Para $\alpha < 0$, note que $|f(x)| > \alpha$ para todo $x \in X$. Desse modo,
  \begin{align*}
    \{x \in X; |f(x)| > \alpha\} = X \in \mathcal F.
  \end{align*}
  Com isso, temos que se $f:X \rightarrow \mathbb C$ é mensurável, segue que $|f|$ também será mensurável.

  \vspace{5mm}
  \noindent
  (Limites de Funções:) Seja $\{f_n(x)\}$ uma sequência de funções mensuráveis a valores complexos
  que converge pontualmente para $f: X \rightarrow \mathbb C$. Vejamos que $f$ é mensurável. 
  
  Sabemos que $n \in \mathbb N$, existem $u_n, v_n: X \rightarrow \mathbb R$ mensuráveis
  tais que $f_n(x) = u_n(x) + iv_n(x)$. Como $\{f_n\}$ converge, 
  \begin{align*}
    \lim f_n(x) = \lim [u_n(x) + iv_n(x)] = [\lim u_n(x)] + i[\lim v_n(x)].
  \end{align*}
  Daqui, já que $\{u_n\}$ e $\{v_n\}$ são sequências de funções mensuráveis que convergem 
  pontualmente para $\Re{f}$ e $\Im{f}$, respectivamente, segue que $\Re{f}$ e $\Im{f}$ são 
  mensuráveis, portanto $\lim f_n(x)= f(x)$ também é mensuráveis.
\end{proof}

\bigline 
\newpage 
\mysubsec{Exercício 6}
\bigline 
\textbf{Seja $f: X \rightarrow Y$ uma função. Mostre que se $\mathcal G$ 
for uma $\sigma$-álgebra de Y, então 
$$\mathcal F= \{f^{-1}(E);E\in\mathcal G\} $$
é uma $\sigma$-álgebra de X.
\begin{proof}[\textbf{Dem:}]
Ora, como $\emptyset \in \mathcal G$, então $f^{-1}(\emptyset) = \emptyset \in \mathcal F$. E, visto que 
  $Y \in \mathcal G$, então $f^{-1}(Y) = X \in \mathcal F$. Daqui, note que para todo $A \in \mathcal F$, 
  existe $B \in \mathcal G$ tal que $f^{-1}(B) = A$. Como $\mathcal G$ é $\sigma$-álgebra de $Y$, temos 
  $B^C \in \mathcal G$, sendo assim, 
  $$f^{-1}(B^C) = \left(f^{-1}(B)\right)^C = A^C \in \mathcal F.$$
  Por fim, 
  note que para toda sequência $\{A_n\}_{n \in \mathbb N} \subset \mathcal F$, existe uma sequência 
  $\{B_n\}_{n \in \mathbb N}$ tal que $f^{-1}(B_n) = A_n$ para todo $n 
  \in \mathbb N$. Como $\bigcup B_n \in \mathcal G$
  \begin{align*}
    f^{-1}\left(\bigcup\limits_{n \in \mathbb N} B_n\right) = \bigcup\limits_{n \in \mathbb N} 
    f^{-1}(B_n) = \bigcup\limits_{n \in \mathbb N} A_n \in \mathcal F.
  \end{align*}
Com isso, temos que $\mathcal F$ é uma $\sigma$-álgebra de $X$, como queríamos. 
\end{proof}
}

\bigline
\mysubsec{Exercício 7}
\bigline
\textbf{Seja $f:X \rightarrow Y$ uma função. Mostre que se $\mathcal F$ 
for uma $\sigma$-álgebra de $X$, então 
$$\mathcal G = \{E \subset Y; f^{-1}(E)\in \mathcal F\}$$ 
é uma $\sigma$-álgebra de Y.
}
\begin{proof}[\textbf{Dem:}]
  Ora, como $X \in \mathcal F$ e $f^{-1}(Y) = X$, então $Y \in \mathcal G$. Mais ainda, 
  $\emptyset \in \mathcal F$ e $f^{-1}(\emptyset) = \emptyset$, portanto $\emptyset \in \mathcal G$.

  Daqui, note que para todo $A \in \mathcal G$, existe $B \in \mathcal F$ tal que $B = f^{-1}(A)$. 
  Visto que $\mathcal F$ é uma $\sigma$-álgebra de $X$, então $B^C \in \mathcal F$. Já que, 
  $$B^C = (f^{-1}(A))^C = f^{-1}\left(A^C\right) \in \mathcal F,$$
  temos $A^C \in \mathcal G$. Por fim, note que para toda sequência $\{A_n\}_{n\in \mathbb N} \subset \mathcal G$,
  existe uma sequência $\{B_n\}_{n \in \mathbb N} \subset \mathcal F$, tal que
  $B_n = f^{-1}(A_n)$ para todo $n \in \mathbb N$. Visto que $\mathcal F$ é uma $\sigma$-álgebra de $X$, temos 
  $\bigcup B_n \in \mathcal F$, isto é,
  \begin{align*}
    \bigcup\limits_{n \in \mathbb N} B_n = \bigcup\limits_{n \in \mathbb N} f^{-1}(A_n) = 
    f^{-1}\left(\bigcup\limits_{n \in \mathbb N} A_n\right) \in \mathcal F,
  \end{align*}
 ou seja, $\bigcup A_n \in \mathcal G$. Com isso, segue que $\mathcal G$ é uma $\sigma$-álgebra de $Y$.
\end{proof}

\bigline
\newpage
\mysubsec{Exercício 8}
\bigline
\textbf{Seja $(X, \mathcal F)$ um espaço mensurável e $f: X \rightarrow Y$ uma função. Seja $\mathcal C$ uma coleção de 
subconjuntos de $Y$ tais que $f^{-1}(A) \in \mathcal F$, para todo $A \in \mathcal C$. Mostre que $f^{-1}(E) \in \mathcal F$, 
para todo $E \in \sigma(\mathcal C)$}.
\begin{proof}[\textbf{Dem:}] Considere a seguinte coleção de subconjuntos de $Y$:
  \begin{align*}
    \mathcal G = \{E\subseteq Y; f^{-1}(E) \in \mathcal F\}.
  \end{align*}
  Do exercício 7, temos que $\mathcal G$ é $\sigma$-álgebra de $Y$. Já que 
  para todo $A \in \mathcal C$ segue que $f^{-1}(A) \in \mathcal F$,
  então $\mathcal C \subseteq \mathcal G$. Da definição de $\sigma$-álgebra gerada, 
  temos $\sigma(\mathcal C) \subseteq \mathcal G$, portanto, para todo 
  $A \in \sigma(\mathcal C)$, segue que $f^{-1}(A) \in \mathcal F$, como queríamos. 
\end{proof}


\bigline
\mysubsec{Exercício 9}
\bigline 
\textbf{Sejam $(X, \mathcal F)$ um espaço mensurável e $f: X \rightarrow \mathbb R$ uma 
função. Mostre que $f$ é mensurável se, e somente se, $f^{-1}(E) \in \mathcal F$, 
para todo $E \in \mathcal B(\mathbb R)$. Conclua que as definições de mensurabilidade
da função $f$ vistas em aula coincidem ao considerar $Y=(\mathbb R, \mathcal B(\mathbb R))$.
}
\begin{proof}[\textbf{Dem:}]Seja  
  $\mathcal C=\{(\alpha, +\infty)\subset \mathbb R; \alpha \in \mathbb R\}$ uma coleção de 
  subconjuntos de $\mathbb R$. 

  \vspace{3mm}
($\implies$) Suponhamos $f$ mensurável, isto é,
  para todo $\alpha \in \mathbb R$, 
  $A_{\alpha} = \{x \in X; f(x) > \alpha\} \in \mathcal F$.  
   Visto que para todo $\alpha \in \mathbb R$, segue que 
  $A_\alpha = f^{-1}((\alpha, +\infty)) \in \mathcal F$, temos do exercício anterior que
  $f^{-1}(E) \in \mathcal F$, para todo $E \in \sigma(\mathcal C)$. Já que 
  $\mathcal C$ gera $\mathcal B(\mathbb R)$, então 
  $\sigma(\mathcal C) = \mathcal B (\mathbb R)$, como queríamos.

  
  ($\impliedby$) Por outro lado, suponhamos que $f^{-1}(E) \in \mathcal F$ para todo $E \in 
  \mathcal B(\mathbb R)$. Como $\mathcal C$ gera $\mathcal B(\mathbb R)$, temos 
  $\mathcal C \subset \mathcal B(\mathbb R)$. Desse modo, para todo $E \in \mathcal C$, 
  segue da hipótese que $f^{-1}(E) \in \mathcal F$. Daqui, note que para todo 
  $E \in \mathcal C$, existe $\alpha \in \mathbb R$ tal que $E = (\alpha, +\infty)$, 
  desse modo, 
  \begin{align*}
    f^{-1}(E) = f^{-1}((\alpha, +\infty)) = \{x\in X; f(x) > \alpha\}.
  \end{align*}
  Portanto, para todo $\alpha \in \mathbb R$, temos 
  $\{x \in X; f(x)>\alpha \} \in \mathcal F$, i.e, $f$ é mensurável, como queríamos.

  Mais ainda, por definição, dizemos que $f: (X,\mathcal F) \rightarrow (Y, \mathcal G)$ é
  mensurável se $f^{-1}(E) \in \mathcal F$, para todo $E \in \mathcal G$. De fato, ao tomar 
  $(Y, \mathcal G)$ como sendo $(\mathbb R, \mathcal B(\mathbb R))$, temos a primeira parte
  da proposição do presente exercício.
\end{proof}

\bigline
\newpage
\mysubsec{Exercício 10}
\bigline 
\textbf{Sejam $(X, \mathcal F)$ um espaço mensurável, $f:X \rightarrow \mathbb R$ 
uma função mensurável e $\varphi: \mathbb R \rightarrow \mathbb R$ uma função contínua. 
Mostre que $\varphi \circ f: X \rightarrow \mathbb R$ é mensurável.}
\begin{proof}[\textbf{Dem:}]
  Ora, como $\varphi$ é contínua, então dado $\alpha \in \mathbb R$ temos $\varphi^{-1}((\alpha, +\infty))$ 
  aberto. Já que todo aberto pode ser escrito como uma união enumerável de elementos de $\mathcal B (\mathbb R)$,
  segue que $$\varphi^{-1}((\alpha, +\infty))\in \mathcal B(\mathbb R).$$
  Daqui, como $f$ é mensurável, para todo $E \in \mathcal B (\mathbb R)$, temos $f^{-1}(E) \in \mathcal F$. Sendo assim, 
 tomando $E = \varphi^{-1}((\alpha, +\infty))$, segue que $f^{-1}(\varphi^{-1}((\alpha, +\infty))) \in 
  \mathcal F$. Dito isso, note que 
  \begin{align*}
    f^{-1}(\varphi^{-1}((\alpha, +\infty))) &= (\varphi \circ f)^{-1}((\alpha, +\infty)),\\
    &=\{x \in X; (\varphi \circ f)(x) > \alpha\} = A_{\alpha}.
  \end{align*}
  Já que para todo $\alpha \in \mathbb R$ temos $A_{\alpha} \in \mathcal F$, então $\varphi \circ f$ é mensurável.
\end{proof}

\bigline
\mysubsec{Exercício 11}
\bigline 
\textbf{Sejam $(X, \mathcal F)$ um espaço mensurável, $f:X \rightarrow \mathbb R$ 
uma função mensurável e $\psi: \mathbb R \rightarrow \mathbb R$ uma função Borel 
mensurável. Mostre que $\psi \circ f: X \rightarrow \mathbb R$ é mensurável.}
\begin{proof}[\textbf{Dem:}]
  Já que $\psi$ é borel mensurável, dado $E \in \mathcal B(\mathbb R)$, 
  temos $\psi^{-1}(E) \in \mathcal B(\mathbb R)$. Visto que $f$ é mensurável,  dado 
  $A \in \mathcal B(\mathbb R)$, segue que $f^{-1}(A) \in \mathcal F$, Desse modo, 
  tomando $A = \psi^{-1}(E)$, temos $f^{-1}(\psi^{-1}(E)) \in \mathcal F$. Note que, 
\begin{align*}
  f^{-1}(\psi^{-1}(E)) = (\psi \circ f)^{-1}(E),
\end{align*}
  ou seja, para todo $E \in \mathcal B(\mathbb R)$, temos $(\psi \circ f)^{-1}(E) \in \mathcal F$, ou seja 
  $\psi \circ f$ é mensurável.
\end{proof}

\bigline
\newpage
\mysubsec{Exercício 12}
\bigline 
\textbf{Sejam $f: X \rightarrow Y$ uma função, $\mathcal F$ uma $\sigma$-álgebra de $X$ e 
$\mathcal C \subset \mathcal P(Y)$. Mostre que 
\begin{align*}
  \sigma(\{f^{-1}(A); A \in \mathcal C\}) = \{f^{-1}(E); E \in \sigma(\mathcal C)\}
\end{align*}
}
\begin{proof}[\textbf{Dem:}] 
  Sejam $D = \{f^{-1}(A) ; A \in \mathcal C\}$ e $\mathcal G = \{f^{-1}(E); E \in \sigma(\mathcal C)\}$.
  Vejamos que $\sigma(D) = \mathcal G$.
  \vspace{3mm}

  ($\subseteq$) Visto que $\mathcal C \subseteq \sigma(\mathcal C)$, então 
  $$D = \{f^{-1}(E); E \in \mathcal C\} \subseteq \{f^{-1}(E); E \in \sigma(\mathcal C)\} = \mathcal G.$$
  Pelo exercício 6, temos que $\mathcal G$ é uma $\sigma$-álgebra de $X$. Visto que $D$ é uma 
  coleção de subconjuntos de $X$, temos $\sigma(D) \subseteq \mathcal G$. 

  ($\supseteq$) Seja $\mathcal A = \{E \subseteq Y; f^{-1}(E) \in \sigma(D)\}$. Do exercício 7, temos que 
  $\mathcal A$ é $\sigma$-álgebra de Y. Com isso, vejamos que $\mathcal C \subset \mathcal A$. 
  Dado $A \in \mathcal C \subseteq Y$, temos $f^{-1}(A) \in D \subseteq \sigma(D)$, pela definição de $\mathcal A$, 
  segue que $A \in \mathcal A$, isto é, $\mathcal C \subseteq \mathcal A$. Das propriedades de $\sigma$-álgebra gerada, 
  temos $\sigma(\mathcal C) \subseteq \mathcal A$. Daqui, note que 
  \begin{align*}
    \mathcal G = \{f^{-1}(E); E \in \sigma(\mathcal C)\} \subseteq 
    \{f^{-1}(E); E \in \mathcal A\} \subseteq \sigma(D)
  \end{align*}
  ou seja $\mathcal G \subseteq \sigma(D)$. Portanto, $\mathcal G = \sigma(D)$, como queríamos.

  \href{https://math.stackexchange.com/questions/1620852/how-to-proof-that-
  f-1-sigma-mathcal-c-subseteq-sigmaf-1-mathcal-c?noredirect=1&lq=1}{(Ver dica recebida)}
\end{proof}

\bigline
\mysubsec{Exercício 13}
\bigline 
\textbf{Sejam $f: X \rightarrow Y$ uma função, $\mathcal F$ uma $\sigma$-álgebra de $X$ e $\mathcal C \subset \mathcal P(Y)$. 
Mostre que $f: (X, \mathcal F) \rightarrow (Y, \sigma(\mathcal C))$ é mensurável se, e somente se, 
para todo $A \in \mathcal C$, tem-se que $f^{-1}(A) \in \mathcal F$.}
\begin{proof}[\textbf{Dem:}]
  A demonstração será dividida em duas partes, são elas:

  ($\implies$) Se $f: X \rightarrow Y$ é mensurável, segue por definição que, para todo $E \in \sigma(\mathcal C)$, 
  temos $f^{-1}(E) \in \mathcal F$. Já que $\mathcal C \subseteq \sigma(\mathcal C)$, então para todo $A \in \mathcal C$, 
  tem-se que $f^{-1}(A) \in \mathcal F$, como queríamos.

  ($\impliedby$) Por outro lado, suponhamos que para todo $A \in \mathcal C$, temos $f^{-1}(A) \in \mathcal F$. Do exercício 
  8, segue que para todo $E \in \sigma(\mathcal C)$, tem-se que $f^{-1}(E) \in \mathcal F$, ou seja, $f$ é mensurável por 
  definição, como queríamos.
\end{proof}
